% Autogenerated translation of 020-test.md by Texpad
% To stop this file being overwritten during the typeset process, please move or remove this header

\documentclass[12pt]{book}
\usepackage{graphicx}
\usepackage{fontspec}
\usepackage[utf8]{inputenc}
\usepackage[a4paper,left=.5in,right=.5in,top=.3in,bottom=0.3in]{geometry}
\setlength\parindent{0pt}
\setlength{\parskip}{\baselineskip}
\setmainfont{Helvetica Neue}
\usepackage{hyperref}
\pagestyle{plain}
\begin{document}

\hrule
layout: post
title: Тесты и упражнения

\section*{slug: epistemic2}

:blue\emph{book: \textbf{Упражнение 1}. Постройте отмеченные эпистемические модели для следующих формул:
1. \textbackslash{}(p wedge neg K}a p\textbackslash{})
2. \textbackslash{}(K\emph{a p wedge neg K}a q\textbackslash{})
3. \textbackslash{}(K\emph{a p wedge neg K}b p\textbackslash{})
4. \textbackslash{}(K\emph{a p wedge neg K}a q wedge K\emph{b q wedge neg K}b p\textbackslash{})

:blue\emph{book: \textbf{Упражнение 2}. Постройте отмеченные эпистемические модели для следующих формул (или докажите, что это невозможно): 
1. \textbackslash{}(hat\{K\}}a p wedge hat\{K\}\emph{a neg p \textbackslash{})
2. \textbackslash{}(K}a p wedge (p to q) wedge neg K\emph{a q \textbackslash{})
3. \textbackslash{}(K}a p wedge K\emph{a ( p to q) wedge neg K}a q \textbackslash{}) 
4. \textbackslash{}(neg K\emph{a p wedge K}b p \textbackslash{})
5. \textbackslash{}(K\emph{a p wedge K}b neg p \textbackslash{})
6. \textbackslash{}( K\emph{a K}b p wedge neg K\emph{a p \textbackslash{})
7. \textbackslash{}( K}a K\emph{b p wedge neg K}b K\emph{a p \textbackslash{})
8. \textbackslash{}( K}b K\emph{a p wedge K}a K\emph{b p wedge neg K}a K\emph{b K}a p \textbackslash{})

:blue\_book: \textbf{Упражнение 3}. Постройте отмеченные эпистемические модели для следующих формул (или докажите, что это невозможно): 

\begin{enumerate}
\item \textbackslash{}( E\emph{\{ab\}p  wedge K}a K\emph{b p wedge neg K}b K\_a p \textbackslash{})
\item \textbackslash{}( E\emph{\{ab\}p  wedge neg K}a K\emph{b p wedge neg K}b K\_a p \textbackslash{})
\item \textbackslash{}( E\textasciicircum{}2\emph{\{ab\}p wedge K}a K\emph{b K}a p wedge neg K\emph{b K}a K\_b p \textbackslash{})
\item \textbackslash{}( E\textasciicircum{}2\emph{\{ab\}p wedge neg K}a K\emph{b K}a p wedge neg K\emph{b K}a K\_b p \textbackslash{})
\end{enumerate}

:blue\_book: \textbf{Упражнение 4}. Постройте эпистемические модели, которые бы описывали следующую ситуацию:

\begin{enumerate}
\item На столе лежит конверт, в котором возможно лежит письмо (пусть \$\$p:=\$\$ $<$$<$в конверте лежит письмо$>$$>$). Ни Аня, ни Борис не знают пуст конверт или нет.
\item Аня и Борис вместе открывают конверт.
\item Аня на глазах Бориса открывает конверт, но не показывает его содержимое Борису.
\item Борис вышел в другую комнату, а затем вернулся. За это время Аня могла посмотреть содержимое конверта.
\item :star: Сначала Борис вышел в другую комнату и вернулся, а потом Аня вышла в другую комнату и вернулась. Каждый из них мог открыть конверт.
\end{enumerate}

:blue\emph{book: \textbf{Упражнение 5}. Найдите формулы которые различают модели в предыдущем упражнении. Важно, что речь идет именно о модели, а не отмеченной модели, то есть, формула должна выполнятся в каждом из миров. Используйте оператор общего знания \$\$C}\{ab\}\$\$. Также удобно использовать следующее сокращение: \$\$K\textasciicircum{}\{?\}\emph{\{i\} varphi:= K}i varphi vee K\_i neg varphi\$\$.

\end{document}
